% !TeX program = lualatex
% halo:begin
% {
%   "version": 1,
%   "publish": true,
%   "course": "mathematical-analysis-i",
%   "section": "2.3",
%   "title": "数学分析（一）2.3：数列极限的计算",
%   "categories": ["study-notes"],
%   "tags": ["mathematics"],
%   "excerpt": "用夹逼定理与四则运算法则计算方根、阶乘、多项式之比、根式及含参数数列的极限。"
% }
% halo:end
% 来源：IMG_1588.HEIC、IMG_1589.HEIC。
\RequirePackage{iftex}
\RequireLuaTeX
\documentclass[UTF8, fontset=fandol, a4paper, 12pt]{ctexart}
\usepackage[margin=2.5cm]{geometry}
\usepackage{amsmath,amssymb,amsthm}
\usepackage{enumitem,fancyhdr}
\usepackage[hidelinks]{hyperref}
\newcommand{\R}{\mathbb R}
\newcommand{\Nplus}{\mathbb N_{+}}
\theoremstyle{definition}
\newtheorem{definition}{定义}
\newtheorem{theorem}{定理}
\renewcommand{\thetheorem}{2.\arabic{theorem}}
\renewcommand{\proofname}{证明}
\setlist[enumerate]{leftmargin=2em,itemsep=0.2em,topsep=0.3em}
\setlength{\headheight}{16pt}
\pagestyle{fancy}
\fancyhf{}
\fancyhead[L]{\small 数学分析（一）}
\fancyhead[R]{\small 2.3\quad 数列极限的计算}
\fancyfoot[C]{\thepage}
\renewcommand{\headrulewidth}{0.3pt}
\raggedbottom
\title{数列极限的计算}
\author{Yuchen Zhang}
\date{整理日期：2026年9月29日}
\makeatletter
\renewcommand{\maketitle}{\begin{center}
{\LARGE\bfseries \@title\par}\vspace{0.6em}
\@author\par\vspace{0.2em}{\small\@date\par}
\end{center}}
\makeatother
\newtheorem{example}{例题}
\begin{document}
\maketitle
\section{方根与阶乘的极限}
\begin{example}
设 $a>0$，证明 $a^{1/n}\to1$；证明 $n^{1/n}\to1$。
\end{example}
\begin{proof}
若 $a>1$，令 $t_n=a^{1/n}-1>0$，由二项式定理，
\[
 a=(1+t_n)^n\ge1+nt_n,\qquad
 0<t_n\le\frac{a-1}{n}\longrightarrow0.
\]
由夹逼定理，$a^{1/n}=1+t_n\to1$。
$a=1$ 时显然成立；若 $0<a<1$，对 $1/a>1$ 应用上式，再取倒数即可。

令 $s_n=n^{1/n}-1\ge0$。当 $n\ge2$ 时，
\[
 n=(1+s_n)^n\ge\binom n2s_n^2,\qquad
 0\le s_n\le\sqrt{\frac2{n-1}}\longrightarrow0.
\]
故 $n^{1/n}\to1$。
\end{proof}

\begin{example}
证明 $\displaystyle\lim_{n\to\infty}\frac1{\sqrt[n]{n!}}=0$。
\end{example}
\begin{proof}
\textbf{方法一。} 已知对每个固定的 $c>0$，有 $c^n/n!\to0$。
任取 $\varepsilon>0$，令 $c=1/\varepsilon$。存在 $N$，使 $n>N$ 时
\[
 \frac{(1/\varepsilon)^n}{n!}<1,
 \qquad 0<\frac1{\sqrt[n]{n!}}<\varepsilon.
\]
故所求极限为 $0$。

\textbf{方法二。} 当 $n\ge2$ 时，$n!$ 中至少有 $n/2$ 个因子不小于 $n/2$，
其余因子均不小于 $1$，因而
\[
 n!\ge\left(\frac n2\right)^{n/2},\qquad
 0<\frac1{\sqrt[n]{n!}}\le\sqrt{\frac2n}\longrightarrow0.
\]
再由夹逼定理得结论。
\end{proof}

\newpage
\section{多项式之比}
\begin{example}
求下列极限：
\[
 \text{（1）}\ \lim_{n\to\infty}\frac{n^3+2n^2}{2n^3+1},\qquad
 \text{（2）}\ \lim_{n\to\infty}\frac{3n^2+2}{2n^3+1},\qquad
 \text{（3）}\ \lim_{n\to\infty}\frac{n^5+2}{2n^4+n^3}.
\]
\end{example}
\begin{proof}[解]
\textbf{（1）} 分子、分母同除以 $n^3$，得
\[
 \frac{n^3+2n^2}{2n^3+1}
 =\frac{1+2/n}{2+1/n^3}\longrightarrow\frac12.
\]
\textbf{（2）} 同除以 $n^3$，得
\[
 \frac{3n^2+2}{2n^3+1}
 =\frac{3/n+2/n^3}{2+1/n^3}\longrightarrow0.
\]
\textbf{（3）} 同除以 $n^4$，得
\[
 \frac{n^5+2}{2n^4+n^3}
 =\frac{n+2/n^4}{2+1/n}\ge\frac n3\longrightarrow+\infty.
\]
所以第三个数列趋于正无穷，不收敛于有限实数。
\end{proof}

\section{根式之差}
\begin{example}
求 $\displaystyle\lim_{n\to\infty}\sqrt n\bigl(\sqrt{n+1}-\sqrt n\bigr)$。
\end{example}
\begin{proof}[解]
将根式之差有理化，得
\[
 \sqrt n\bigl(\sqrt{n+1}-\sqrt n\bigr)
 =\frac{\sqrt n}{\sqrt{n+1}+\sqrt n}
 =\frac1{\sqrt{1+1/n}+1}\longrightarrow\frac12.
\]
\end{proof}

\newpage
\section{含参数的极限}
\begin{example}
设 $a\in\R\setminus\{-1\}$，求
$\displaystyle\lim_{n\to\infty}\frac{a^n}{a^n+1}$。
\end{example}
\begin{proof}[解]
分三种情形讨论。
\begin{enumerate}
\item 若 $a=1$，则各项均为 $1/2$，故极限为 $1/2$。
\item 若 $|a|<1$，由 $a^n\to0$，得 $a^n/(a^n+1)\to0$。
\item 若 $|a|>1$，则 $(1/a)^n\to0$，从而
\[
 \frac{a^n}{a^n+1}=\frac1{1+(1/a)^n}\longrightarrow1.
\]
\end{enumerate}
因此
\[
 \lim_{n\to\infty}\frac{a^n}{a^n+1}
 =\begin{cases}
 0,&|a|<1,\\
 \tfrac12,&a=1,\\
 1,&|a|>1.
 \end{cases}
\]
$a=-1$ 时奇数项的分母为零，不能构成题设的数列。
\end{proof}

\end{document}
